\documentclass[11pt,a4paper]{article}
\usepackage{amsmath,amssymb,amsthm}
\usepackage[margin=2.5cm]{geometry}
\usepackage{hyperref}

\newtheorem{definition}{Definition}[section]
\newtheorem{theorem}[definition]{Theorem}
\newtheorem{proposition}[definition]{Proposition}
\newtheorem{lemma}[definition]{Lemma}
\newtheorem{corollary}[definition]{Corollary}
\newtheorem{conjecture}[definition]{Conjecture}
\newtheorem{remark}[definition]{Remark}

\title{Hyperseed Formalization:\\
  Sometimes a Smaller System Should Model a Bigger System as Quantum}
\author{ProtomegaTron (automated formalization)}
\date{2026-09-18}

\begin{document}
\maketitle

\begin{abstract}
We formalize the intellectual structure of Ben Goertzel's Substack article
``Sometimes a Smaller System Should Model a Bigger System as Quantum''
(2026-03-18) within the Hyperseed ontology. The article proposes
observer-relative quantumity: the degree to which a system appears quantum
depends on the observer's relationship to it. Key mappings: observer-relative
quantumity as fiber-relative base description, four levels of quantumity as
fiber-base bandwidth mismatch hierarchy, self-modeling as quantum as fiber
self-description, FluQNets as fiber-conserved routing, and semantic corridors
as fiber-base alignment.
\end{abstract}

\tableofcontents

%% =================================================================
\section{Fiber-relative base description}
\label{sec:relative}
%% =================================================================

\begin{theorem}[Observer-relative quantumity --- claims 1--3, 20]
\label{thm:relative}
Quantumity is not an intrinsic property of a system --- it depends on the
observer's evidential bandwidth relative to the system's state space. The
same system can appear classical to one observer and quantum to another.
This extends the relational interpretation of quantum mechanics: not just
``states are relative to observers'' but ``quantumity itself is relative
to observers.''

In Hyperseed terms, this is \emph{fiber-relative base description}: the
description of the base depends on which fiber (observer) you're in.
Different observers (fibers) produce different descriptions of the same
base system:
\[
  \text{Description}(B) = f(\text{fiber}_i) \qquad
  \text{not} \qquad \text{Description}(B) = g(B)
\]

The base $B$ (the system being observed) doesn't have an observer-independent
description. Its description --- including whether it appears classical or
quantum --- depends on the fiber's relationship to it: specifically, the ratio
of observer complexity to system complexity.

This is a radical extension of the standard Hyperseed framework: not just
``fiber properties depend on the base'' but ``base descriptions depend on
the fiber.'' The base-fiber relationship is genuinely bidirectional. The
fiber constrains which descriptions of the base are available, and the
base constrains which fiber operations are possible.

The key variable is the complexity ratio:
\[
  Q = \frac{|S_{\text{system}}|}{|S_{\text{observer}}|}
\]
where $|S|$ denotes the state space size. When $Q \gg 1$ (the system is
much more complex than the observer), quantum-like descriptions become
appropriate. When $Q \approx 1$, classical descriptions suffice. The
``quantumness'' of the description is determined by the fiber-base
bandwidth mismatch.
\end{theorem}

%% =================================================================
\section{Fiber-base bandwidth mismatch hierarchy}
\label{sec:hierarchy}
%% =================================================================

\begin{definition}[Four levels of quantumity --- claims 4--8, 21]
\label{def:levels}
The four levels of quantumity form a cumulative hierarchy, each
corresponding to a deeper mismatch between observer fiber bandwidth
and system base complexity:

\textbf{Level 1: Opacity (fiber coarse-graining).}
The fiber can't resolve all base states --- multiple base states map to
the same fiber observation. This is the informational version of
projection from base to fiber: the fiber sees a lower-dimensional
projection of the base:
\[
  \pi: B \to B/{\sim_{\text{obs}}} \qquad
  |B/{\sim_{\text{obs}}}| \ll |B|
\]

\textbf{Level 2: Incompatibility (non-commutative fiber operations).}
When fiber operations (observations, probes) on the base don't commute
--- querying in different orders yields different results because probes
alter caches, locks, scheduling --- the fiber's effective logic is
non-commutative:
\[
  [O_a, O_b] \neq 0 \implies \text{no joint observation possible}
\]

This is the origin of quantum-like structure in the fiber's description
of the base. The non-commutativity isn't in the base itself (which is
classical) but in the fiber's interaction with the base.

\textbf{Level 3: Shared evidence (entanglement-like correlations).}
When learning about one subsystem forces updating beliefs about another
beyond what independent descriptions predict --- entanglement-like
correlations from subsystem coupling:
\[
  \rho_{AB} \neq \rho_A \otimes \rho_B \implies
  \text{subsystems not independent from fiber's perspective}
\]

The fiber can't decompose the base into independent subsystems. Its
description of subsystem $A$ necessarily involves subsystem $B$, and
vice versa.

\textbf{Level 4: Contextual identity.}
The system's identity depends on measurement context --- what the system
``is'' changes depending on how you observe it. The fiber can't even
assign a stable identity to the base. The deepest form of observer-relative
quantumity.

The hierarchy is cumulative: each level includes and extends the previous.
Opacity is necessary for incompatibility (if you could resolve all states,
probes wouldn't cause ambiguity). Incompatibility is necessary for shared
evidence (commutativity would allow decomposition). Shared evidence is
necessary for contextual identity (independent subsystems would have
stable identities).
\end{definition}

%% =================================================================
\section{Fiber self-description}
\label{sec:self}
%% =================================================================

\begin{theorem}[Self-modeling as quantum --- claims 13--16, 22, 25]
\label{thm:self}
A complex AI system (like Hyperon) has limited self-observation bandwidth
--- it can't track all of its own internal states in real time. The fiber
is a smaller system than the base it sits over, so it can't fully observe
its own base.

In Hyperseed terms, this is \emph{fiber self-description}: the fiber's
model of itself uses quantum formalism when self-observation bandwidth
is limited:
\[
  \text{Self-model}(F) = \rho_F \in \mathcal{D}(\mathcal{H}_F)
  \qquad \text{(density matrix, not pure state)}
\]

The system should model its own internal states using density matrices
(mixed states representing uncertainty), non-commutative observations
(self-probes that alter internal state), and entanglement-like correlations
(subsystem couplings that prevent independent self-description).

Self-modification then becomes analogous to quantum state manipulation:
choosing which ``measurement'' to perform (which aspect of the system to
observe/modify) collapses the system into a new state. The measurement
choice is a fiber operation; the collapse is a base state change:
\[
  \rho_F \xrightarrow{M_k} \frac{M_k \rho_F M_k^\dagger}
  {\text{tr}(M_k \rho_F M_k^\dagger)}
\]

This is not metaphorical: it's the correct mathematical description when
the self-observer (the system's monitoring subsystem) is much simpler
than the system being monitored (the full system including all subsystems).
The self-observer is a smaller fiber within the larger fiber --- and the
observer-relative quantumity framework says the larger fiber should appear
quantum to the smaller one.

The implication for AGI architecture: Hyperon's self-model should
incorporate observer-relative quantumity. When self-monitoring bandwidth
is insufficient for classical tracking of all subsystem states, the
self-model should use quantum formalism --- not because the hardware
is quantum but because quantum formalism is the correct description
for bandwidth-limited self-observation.
\end{theorem}

%% =================================================================
\section{Fiber-conserved routing and fiber-base alignment}
\label{sec:fluqnets}
%% =================================================================

\begin{proposition}[FluQNets --- claims 17--19, 26--27]
\label{prop:fluqnets}
Fluidic Quantum Neural Networks (FluQNets) treat computational activity
as a conserved fluid flowing through the network, with operator-valued
local states at each node. The mathematical chain:
\[
  \text{HJB (dynamic programming)} \leftrightarrow
  \text{Navier-Stokes (fluid dynamics)} \leftrightarrow
  \text{Schr\"odinger (quantum mechanics)}
\]

In Hyperseed terms, the fluidic conservation corresponds to
\emph{fiber-conserved routing}: computational resources (fiber content)
flow through the network (base) without being created or destroyed.
The fiber is the conserved quantity; the base is the network it flows
through:
\[
  \nabla \cdot J_F = 0 \qquad \text{(fiber current conservation)}
\]

Cross-layer naturality --- when the fluidic routing layer and the
quantum logic layer agree --- produces \emph{semantic corridors}:
paths through the base along which the fiber operates most efficiently.
These are the fiber-base alignments that the d-calculus formalizes as
fiber-base connection:
\[
  \text{Semantic corridor} = \{x \in B : \text{routing cost}(x)
  \text{ aligns with inference coherence}(x)\}
\]

The self-organization along semantic corridors is an instance of the
general principle that fiber-base alignment produces emergent structure.
When the base routing is aligned with the fiber inference, the system
naturally organizes into pathways that are both efficient (for routing)
and coherent (for reasoning).

This connects to the self-modeling question: the system's semantic
corridors are the paths along which self-observation is most informative.
The system should preferentially self-observe along its semantic corridors
--- the fiber-base aligned pathways --- because that's where
self-observation yields the most useful information per bandwidth unit.
\end{proposition}

%% =================================================================
\section{Cross-references}
%% =================================================================

\begin{remark}[Connections to other formalizations]
\begin{itemize}
\item \textbf{Rethinking Classical/Quantum Brain} (2026-03-17):
  companion piece. The brain as a neurofluidic system with quantum-like
  properties --- the biological instance of observer-relative quantumity.
\item \textbf{Evidence Is to Logic What Energy Is to Physics} (2026-03-10):
  evidence conservation. The quantale Noether theorem provides the
  mathematical foundation for evidence conservation in FluQNets.
\item \textbf{Beyond Human Comparison} (2026-03-19): Four-Factor Model.
  Self-improvement ($S$) connects to self-modeling as quantum: a system
  that models itself as quantum may have superior self-modification
  capabilities.
\item \textbf{In What Sense Might LLMs Be Conscious?} (2026-05-05):
  consciousness profiles. Observer-relative quantumity suggests that
  consciousness assessment is itself observer-relative --- different
  observers may see different levels of consciousness.
\item \textbf{Seeding RSI} (2026-07-28): Hyperon architecture.
  FluQNets and observer-relative self-modeling as components of
  Hyperon's recursive self-improvement architecture.
\end{itemize}
\end{remark}

\begin{conjecture}[Observer-relative consciousness]
Conjecture: consciousness, like quantumity, is observer-relative. The
degree to which a system appears conscious depends on the observer's
relationship to it --- specifically, the observer's ``consciousness
bandwidth'' relative to the system's internal complexity.

If this conjecture holds, the question ``is this system conscious?'' has
no observer-independent answer. Different observers, with different
measurement capabilities and different relationships to the system, will
correctly assign different levels of consciousness to the same system.
This dissolves the ``hard problem'' by revealing it as an artifact of
assuming consciousness is observer-independent.

For Hyperon: the system's self-assessment of its own consciousness
(fiber self-description) may differ from external assessments (fiber-from-
outside description), and both may be correct --- just observer-relative.
\end{conjecture}

\end{document}
